% 2.3 — brief en documento/investigacion/2-3-ingenieria-de-datos.md
\subsection{INGENIERÍA DE DATOS}
\subsubsection{PROCESOS ETL Y ELT}

Los procesos de extracción, transformación y carga (ETL, por \textit{extract, transform, load}) son el mecanismo clásico para alimentar un almacén de datos desde los sistemas operacionales. \textcite{Vassiliadis2009} los define como los procesos de software que facilitan la carga original y el refresco periódico del contenido del almacén, y revisa los problemas propios de cada una de las tres etapas y los que afectan al proceso completo. El destino es el almacén de datos que \textcite{Inmon2005} caracteriza como una colección de datos orientada a temas, integrada, variable en el tiempo y no volátil, en apoyo de la toma de decisiones; la integración y la orientación a temas son, precisamente, el trabajo que ETL realiza sobre datos que en origen están organizados por aplicación y por transacción.

Desde la práctica del modelado dimensional, \textcite{Kimball2013} advierten que el sistema ETL consume una parte desproporcionada del tiempo y del esfuerzo necesarios para construir un entorno de almacén de datos e inteligencia de negocios, y lo descomponen en 34 subsistemas agrupados en cuatro áreas: extraer los datos de las fuentes, limpiarlos y conformarlos, entregarlos preparados para la presentación y gestionar el entorno ETL\@. Dos de sus aportes interesan a este proyecto. El primero es el conformado (\textit{conforming}), orientado a que atributos equivalentes de fuentes distintas se representen de manera uniforme. El segundo es la distinción entre la carga histórica única, que puebla el almacén con todo el pasado disponible, y el procesamiento incremental, que lo actualiza periódicamente con los registros nuevos.

La variante ELT (\textit{extract, load, transform}) invierte las dos últimas etapas: los datos se cargan crudos en el destino y la transformación se ejecuta allí, con el motor de consulta del propio almacén o lago. \textcite{ReisHousley2022} ubican ambos patrones en la etapa de transformación de un ciclo de vida de la ingeniería de datos que va de la generación a la ingesta, el almacenamiento, la transformación y el servicio; explican que ETL surgió cuando el cómputo era un recurso escaso y convenía transformar antes de cargar, que ELT es el patrón propio de los lagos de datos y \textit{lakehouses} modernos, y recomiendan elegir según el caso en lugar de adherir a uno por principio. \textcite{Armbrust2021} muestran el costo de encadenar ambos patrones: en la arquitectura de dos niveles habitual en las empresas los datos primero se someten a ETL hacia el lago y luego a ELT hacia el almacén, y cada paso adicional agrega demoras, complejidad y nuevas oportunidades de introducir errores que reducen la calidad de los datos.

El pipeline de datos de este proyecto es un caso de ELT con fines exclusivamente analíticos. Las fuentes son la base MySQL del portal anterior, la base PostgreSQL del ERP nuevo y las planillas Excel de control de proyectos que la empresa lleva fuera del sistema; sus datos se copian sin modificaciones a una capa cruda y se transforman después, en las capas que describe la subsección siguiente. La distinción de Kimball y Ross se traduce en dos modos de operación: una carga histórica del portal anterior y de las planillas, que produce el conjunto de entrenamiento, y una carga incremental desde el ERP nuevo, que alimenta la inferencia sobre los proyectos activos. En ningún caso el pipeline escribe de vuelta en el ERP ni migra el histórico al esquema nuevo; esa migración queda fuera del alcance del trabajo.

\subsubsection{ARQUITECTURA MEDALLION (BRONZE, SILVER, GOLD)}

El principio de organizar los datos en capas de calidad creciente tiene su fundamento académico en la propuesta de \textit{lakehouse} de \textcite{Armbrust2021}. Los autores parten de un diagnóstico: los lagos de datos de segunda generación, que aplican el esquema en el momento de la lectura, trasladaron aguas abajo el problema de la calidad y la gobernanza, y combinarlos con un almacén separado produce datos desactualizados, poca fiabilidad y un alto costo total. En respuesta proponen una plataforma única que almacena los datos en formatos abiertos de acceso directo, como Parquet, ofrece soporte de primera clase al aprendizaje automático y a la ciencia de datos, y alcanza un rendimiento comparable al de un almacén. Dos rasgos de esa propuesta importan aquí: un \textit{lakehouse} debe poder guardar los datos crudos, como un lago, y al mismo tiempo soportar procesos ETL/ELT que los curen para mejorar su calidad con vistas al análisis; y su capa de metadatos es el lugar natural para aplicar controles de calidad, como la imposición de esquema y las restricciones que rechazan o ponen en cuarentena los registros inválidos.

La nomenclatura con la que hoy se designa ese curado por etapas no proviene de ese artículo, sino de la documentación de Databricks, es decir, de literatura de proveedor. \textcite{DatabricksMedallion} define la arquitectura \textit{medallion}, también llamada de múltiples saltos (\textit{multi-hop}), como un patrón de diseño para organizar lógicamente los datos en una serie de capas que denotan la calidad de lo almacenado. La capa Bronze recibe los datos crudos de las fuentes sin validación y conserva su estado original; la capa Silver contiene al menos una representación validada y no agregada de cada registro, y es donde se realizan la limpieza, la deduplicación y la normalización; la capa Gold consiste en datos agregados y ajustados a los requisitos de análisis y reporte. Al hacer progresar los datos por las capas, su calidad y fiabilidad mejoran de manera incremental hasta hacerlos aptos para inteligencia de negocios y aprendizaje automático.

El proyecto lo adopta sin la plataforma del proveedor, en línea con el argumento de \textcite{Armbrust2021} a favor de los formatos abiertos: se prevé materializar las capas como archivos en el servidor de la empresa y ejecutar las transformaciones con dos herramientas ligeras que se describen en la sección 2.7, DuckDB para Silver y Polars para Gold. El cuadro~\ref{tab:medallion-proyecto} resume la correspondencia. Bronze reúne copias crudas de las tres fuentes: las tablas de MySQL del portal anterior, las tablas de PostgreSQL del ERP nuevo y las planillas Excel. Silver es la capa en la que el conformado de \textcite{Kimball2013} se convierte en un requisito del modelo predictivo: los dos esquemas relacionales describen los mismos objetos de negocio (centros de costo, certificaciones, órdenes de compra, presupuestos) con nombres y estructuras distintos, y DuckDB los mapea a un conjunto único de variables con las mismas unidades y dominios, de modo que un modelo entrenado sobre el histórico pueda predecir sobre los proyectos activos. Gold construye con Polars las tablas de características por proyecto a una fecha de corte, único insumo de los modelos. El diseño detallado de cada capa se desarrolla en el Capítulo III\@.

\begin{table}[H]
  \centering
  \caption{Capas de la arquitectura \textit{medallion} y su correspondencia en el proyecto}\label{tab:medallion-proyecto}
  \begin{tblr}{colspec={X[0.38,l] X[1.21,l] X[1.41,l]}}
    \textbf{Capa} & \textbf{Contenido según el patrón} & \textbf{En el proyecto} \\
    Bronze & Datos crudos de las fuentes, sin validación, en su estado original & Copias de las tablas MySQL del portal anterior, de las tablas PostgreSQL del ERP nuevo y de las planillas Excel de control de proyectos \\
    Silver & Al menos una representación validada y no agregada de cada registro; limpieza, deduplicación y normalización & Esquema unificado en DuckDB: los dos esquemas relacionales y las planillas mapeados a las mismas variables, con tipos y dominios conformados y registros deduplicados \\
    Gold & Datos agregados y ajustados a los requisitos de análisis y reporte & Tablas de características por proyecto construidas con Polars; entrada del entrenamiento (histórico) y de la inferencia (proyectos activos) \\
  \end{tblr}
  \fuente{Elaboración propia, 2026, en base a \textcite{DatabricksMedallion} y \textcite{Armbrust2021}.}
\end{table}

\subsubsection{CALIDAD DE DATOS}

La calidad de los datos no se reduce a su exactitud. \textcite{WangStrong1996} observan que los esfuerzos de mejora tienden a concentrarse en la exactitud, mientras que los consumidores de datos manejan una conceptualización mucho más amplia, y construyen un marco jerárquico que organiza las dimensiones de calidad en cuatro categorías: intrínseca, las propiedades que los datos tienen por sí mismos; contextual, que exige considerarlos en el contexto de la tarea en la que se usan; representacional, referida a la forma en que se presentan; y de accesibilidad. Datos de alta calidad son, en esa formulación, intrínsecamente buenos, apropiados para la tarea, claramente representados y accesibles. De la categoría contextual se sigue que la calidad de un conjunto de datos no se juzga en abstracto, sino respecto de la tarea a la que se destina.

Las dimensiones concretas más usadas son la exactitud, la completitud, la consistencia y la oportunidad, a las que \textcite{BatiniScannapieco2016} dedican un capítulo de su tratado junto con otras dimensiones temporales; el cuerpo de conocimiento de \textcite{DAMA2017} agrega, entre otras, la unicidad y la validez. \textcite{Batini2009} muestran, además, que la literatura ofrece una amplia gama de metodologías de evaluación y mejora, que difieren en fases, estrategias y tipos de datos y de sistemas. El cuadro~\ref{tab:dimensiones-calidad} fija las definiciones operativas que este trabajo adopta para seis dimensiones aplicables a un histórico relacional, y la verificación con que cada una se mide.

\begin{table}[H]
  \centering
  \caption{Dimensiones de calidad de datos adoptadas en el proyecto}\label{tab:dimensiones-calidad}
  \begin{tblr}{colspec={X[0.48,l] X[1.21,l] X[1.31,l]}}
    \textbf{Dimensión} & \textbf{Definición operativa} & \textbf{Verificación prevista} \\
    Exactitud & El valor registrado coincide con el valor real del objeto que describe & Contraste de montos y fechas del sistema contra las planillas de control y documentos de respaldo, sobre una muestra \\
    Completitud & Los valores requeridos para la tarea están presentes & Proporción de proyectos con presupuesto, fechas y certificaciones registrados; patrón de valores nulos por variable \\
    Consistencia & Los valores no se contradicen entre registros ni entre tablas & Sumas de detalle frente a totales de cabecera; estados de flujo compatibles con las fechas \\
    Oportunidad & El dato está disponible cuando la tarea lo necesita & Rezago entre la fecha del evento y la fecha de registro en el sistema \\
    Unicidad & Cada objeto de negocio aparece una sola vez & Duplicados de proyectos, clientes y proveedores por clave natural \\
    Validez & Los valores pertenecen al dominio definido & Valores fuera de rango o de catálogo; tipos incorrectos en campos de texto libre \\
  \end{tblr}
  \fuente{Elaboración propia, 2026, en base a \textcite{WangStrong1996}, \textcite{BatiniScannapieco2016} y \textcite{DAMA2017}.}
\end{table}

La relación entre esas dimensiones y el desempeño de un modelo de aprendizaje automático cuenta con evidencia empírica reciente. \textcite{Mohammed2025} sostienen que datos de entrenamiento incompletos, erróneos o inapropiados conducen a modelos poco confiables, y estudian de forma sistemática el efecto de seis dimensiones de calidad sobre el desempeño de 19 algoritmos de clasificación, regresión y agrupamiento sobre datos tabulares, contaminando los datos de entrenamiento, los de prueba o ambos; un diseño pertinente para un proyecto cuyos datos son tabulares y cuyos modelos se entrenan sobre un histórico y se aplican sobre datos nuevos. \textcite{Sambasivan2021} aportan la dimensión organizacional: a partir de entrevistas con 53 practicantes de inteligencia artificial en contextos de alto riesgo, describen las cascadas de datos, efectos negativos que se propagan aguas abajo a partir de problemas en los datos, y las encuentran omnipresentes, invisibles, de efecto retardado y con frecuencia evitables, fruto de prácticas que subvaloran los datos frente al modelado. \textcite{Gudivada2017} agregan que, para el aprendizaje automático, la calidad va más allá de la limpieza y las transformaciones y requiere gobernar el ciclo de vida completo de los datos.

De esa literatura el proyecto toma dos decisiones. La primera es que la exploración del histórico del Capítulo III usa las dimensiones del cuadro~\ref{tab:dimensiones-calidad} como instrumento y las juzga respecto de la tarea predictiva concreta; qué grado de calidad tiene ese histórico en cada dimensión es un resultado de esa exploración y no se anticipa aquí. La segunda es que la mejora se aborda por dos vías: la limpieza en la capa Silver, que corrige los datos, y la corrección de los defectos del ERP nuevo, que corrige el proceso del que saldrán los datos de inferencia, de modo que los registros sobre los que se realiza la inferencia nazcan con mejor calidad que los históricos con los que se entrena.
